\chapter{Small Bowel Obstruction}
\index{Small Bowel Obstruction}
\label{sec:Small Bowel Obstruction}

\section{Overview}
Small bowel obstruction (SBO) is a partial or complete blockage of the small intestine.
\section{Causes}
This condition results from disruption of normal small bowel obstruction function.

\begin{itemize}[noitemsep]
  \item The most common cause is adhesions from prior abdominal surgery (60\% of cases)
  \item Other causes include hernias (inguinal, incisional), malignancy, volvulus, intussusception (in children), inflammatory bowel disease strictures, and gallstone ileus. This condition happens because of mechanical blockage of intestinal contents.
\end{itemize}
\section{Risk Factors}


\begin{itemize}[noitemsep]
  \item Genetic predisposition and family history
  \item Advancing age
  \item Lifestyle factors (diet, exercise, smoking, alcohol)
  \item Environmental exposures
  \item Underlying medical conditions
  \item Medication use
\end{itemize}
\section{Signs and Symptoms}

\subsection{Common symptoms}
\begin{itemize}[noitemsep]
  \item You may experience crampy abdominal pain (periumbilical), nausea, vomiting (bilious early, feculent late), abdominal distension, and obstipation. 
  \end{itemize}

\subsection{Emergency warning signs}
\begin{itemize}[noitemsep]
  \item Severe or worsening symptoms of small bowel obstruction require immediate medical evaluation
\end{itemize}
\section{Progression and Stages}
The condition may progress through different stages of severity over time. Early stages often involve mild or subtle symptoms that may be overlooked. Moderate stages involve more noticeable symptoms affecting daily function. Advanced stages may involve significant impairment and complications.
\section{Complications}
\begin{itemize}[noitemsep]
  \item Disease progression leading to worsening symptoms
  \item Secondary infections or injuries
  \item Side effects of treatment
  \item Reduced quality of life
  \item Emotional and psychological impact
\end{itemize}
\section{Diagnosis}
\begin{itemize}[noitemsep]
  \item Doctors diagnose this based on abdominal X-ray showing air-fluid levels and dilated proximal bowel loops with collapsed distal bowel
  \item CT abdomen/pelvis is more sensitive and identifies the transition point and cause.
  \item Comprehensive medical history and physical examination
  \item Laboratory tests to assess organ function
  \item Imaging studies to visualize affected structures
  \item Specialized testing depending on the affected system
  \item Consultation with relevant specialists
\end{itemize}
\section{Prevention}
\begin{itemize}[noitemsep]
  \item Maintain a healthy lifestyle with balanced diet and exercise
  \item Avoid known risk factors
  \item Attend regular medical check-ups for early detection
  \item Follow recommended screening guidelines
\end{itemize}
\section{Treatment Options}
\begin{itemize}[noitemsep]
  \item Treatment options include nasogastric decompression, IV fluid resuscitation, electrolyte correction, and close monitoring
  \item Complete obstruction or signs of strangulation (fever, tachycardia, peritoneal signs, metabolic acidosis) require emergent surgical exploration.
  \item Medications to manage symptoms and address underlying mechanisms
  \item Lifestyle modifications and self-care measures
  \item Physical therapy and rehabilitation
  \item Surgical intervention when indicated
  \item Regular monitoring and follow-up care
\end{itemize}
\section{Living with It}
\begin{itemize}[noitemsep]
  \item Follow your treatment plan as prescribed
  \item Maintain regular follow-up with your healthcare provider
  \item Make necessary lifestyle adjustments
  \item Seek support from family, friends, or support groups
  \item Stay informed about your condition
\end{itemize}
\section{Outlook}
Your outlook depends on cause and presence of strangulation. The outlook varies depending on the specific condition, severity at diagnosis, and response to treatment. Early intervention generally leads to better outcomes. Many conditions can be effectively managed with appropriate treatment. Regular follow-up care is important for monitoring and treatment adjustment.
